\documentclass[11pt,a4paper]{article}
\usepackage{amsmath,amssymb,amsthm}
\usepackage[margin=2.5cm]{geometry}
\usepackage{hyperref}

\newtheorem{definition}{Definition}[section]
\newtheorem{theorem}[definition]{Theorem}
\newtheorem{proposition}[definition]{Proposition}
\newtheorem{lemma}[definition]{Lemma}
\newtheorem{corollary}[definition]{Corollary}
\newtheorem{conjecture}[definition]{Conjecture}
\newtheorem{remark}[definition]{Remark}

\title{Hyperseed Formalization:\\
  Playing Around with the ARC-AGI-3}
\author{ProtomegaTron (automated formalization)}
\date{2026-09-18}

\begin{document}
\maketitle

\begin{abstract}
We formalize the intellectual structure of Ben Goertzel's Substack article
``Playing Around with the ARC-AGI-3'' (2026-05-08) within the Hyperseed
ontology. The article reports SingularityNET's 33\% human-normalized score
on ARC-AGI-3 using executable world models with LLM-driven coding and
formal verification. The five-stage loop (observe-model-verify-plan-execute)
maps onto a fiber inspection loop; the executable world model is an explicit
fiber representation; verification is a cocycle consistency check;
informative failures are fiber defect localizations; and the LLM + verifier
architecture is generate-and-verify on fiber hypotheses.
\end{abstract}

\tableofcontents

%% =================================================================
\section{The fiber inspection loop}
\label{sec:loop}
%% =================================================================

\begin{theorem}[Five-stage loop --- claims 6--10, 21]
\label{thm:loop}
The system's core architecture is a five-stage loop:
\begin{enumerate}
\item \textbf{Observe:} receive environment state.
\item \textbf{Model:} LLM writes/edits Python code representing
  state, transitions, goals.
\item \textbf{Verify:} formal verifier replays all observations
  against the model.
\item \textbf{Plan:} search within the verified model for a
  path to the goal.
\item \textbf{Execute:} commit actions from the plan, checking
  predictions against reality step by step.
\end{enumerate}

In Hyperseed terms, this is a \emph{fiber inspection loop}:
\begin{itemize}
\item \textbf{Observe} = inspect the environment fiber (read the
  current section).
\item \textbf{Model} = update the model fiber (revise the agent's
  representation of the environment's fiber structure).
\item \textbf{Verify} = check cocycle consistency (do the model's
  transition functions compose consistently with all observed
  transitions?).
\item \textbf{Plan} = compute a path in the model fiber (find a
  section that reaches the goal).
\item \textbf{Execute} = act in the environment fiber and compare
  (test whether the model fiber matches the environment fiber).
\end{itemize}

The key design insight: most reasoning happens in the \emph{model
fiber} (cheap, reversible simulation), not the \emph{environment
fiber} (expensive, irreversible action). This separation of fiber
exploration from fiber commitment is what makes the system efficient:
the agent explores many possible fiber paths in simulation before
committing to one in reality.
\end{theorem}

%% =================================================================
\section{Executable world model as explicit fiber representation}
\label{sec:model}
%% =================================================================

\begin{definition}[Code as fiber --- claims 6, 22]
\label{def:model}
The executable world model is a Python codebase containing:
\begin{itemize}
\item State representation (the fiber's state space).
\item Transition dynamics (the fiber's transition functions).
\item Goal conditions (target sections of the fiber).
\item A planner (path-finding in the fiber).
\end{itemize}

In Hyperseed terms, this is an \emph{explicit fiber representation}:
the code \emph{is} the fiber structure, written in a form that can
be inspected, verified, revised, and executed. This explicitness is
what enables verification --- you can't verify a black-box neural
network's world model against observations the way you can verify
an explicit codebase.

The choice of code as substrate is significant: code is symbolic,
inspectable, compositional, and verifiable. An implicit world model
(hidden in neural network weights) might capture the same dynamics
but can't be verified component by component. The explicit fiber
representation enables \emph{targeted fiber repair}: when verification
fails, you know which function to revise, not just that the overall
model is wrong.
\end{definition}

%% =================================================================
\section{Verification as cocycle check}
\label{sec:verify}
%% =================================================================

\begin{proposition}[Transition-by-transition consistency --- claims 8, 16, 23]
\label{prop:verify}
The verifier replays all recorded observations against the model and
confirms rules transition by transition. A model revision is accepted
only if it's consistent with \emph{every} observed transition.

In Hyperseed terms, this is a \emph{cocycle consistency check}: the
model's transition functions must compose consistently with the
observed transitions:
\[
  \forall (s_t, a_t, s_{t+1}) \in \text{History}: \quad
  T_{\text{model}}(s_t, a_t) = s_{t+1}
\]

If any observed transition violates the model's predicted transition,
the cocycle fails at that point. The model is rejected until it
satisfies the full cocycle condition.

This is a strong form of verification: not just checking the model
against the most recent observation but against \emph{all} observations.
The cocycle condition is cumulative --- each new observation adds a
constraint, and the model must satisfy all of them simultaneously.
This prevents the model from ``forgetting'' earlier observations when
accommodating new ones.
\end{proposition}

%% =================================================================
\section{Informative failure as fiber defect localization}
\label{sec:failure}
%% =================================================================

\begin{theorem}[Targeted repair --- claims 14--17, 24]
\label{thm:failure}
When a verified plan fails in execution (the environment produces a
state the model didn't predict), the failure is \emph{informative}:
it points at exactly which part of the model was wrong.

In Hyperseed terms, this is \emph{fiber defect localization}: the
failure identifies which transition function is incorrect. The defect
is localized to a specific fiber component, enabling targeted repair
rather than global model replacement.

The cycle of hypothesis formation, verification, execution, and
mismatch-driven revision is a systematic fiber repair process:
\begin{enumerate}
\item Form a hypothesis about a transition function (e.g., ``key
  cycles every K moves'').
\item Verify against all known observations (cocycle check).
\item Execute until mismatch.
\item Localize the defect (``the shape change is tile-triggered,
  not move-triggered'').
\item Revise the specific transition function.
\item Re-verify against all observations.
\end{enumerate}

Each iteration narrows the space of possible fiber defects. The
process converges when the model's fiber structure matches the
environment's fiber structure on all observed transitions.
\end{theorem}

%% =================================================================
\section{Generate-and-verify on fiber hypotheses}
\label{sec:genverify}
%% =================================================================

\begin{proposition}[LLM + verifier --- claims 10, 26]
\label{prop:genverify}
The architecture is explicitly neural-symbolic: the LLM proposes
(generates candidate fiber structures), the verifier disposes
(checks cocycle consistency). Code is the shared substrate.

In Hyperseed terms, this is \emph{generate-and-verify on fiber
hypotheses}: the same pattern as un-jailbreakable AI (the LLM
generates, a separate system verifies), applied to world-modeling
rather than safety.

The generate-and-verify pattern is fiber-appropriate because:
\begin{itemize}
\item Generation (LLM) operates at the fiber level: proposing
  structural hypotheses about how the world works.
\item Verification operates at the section level: checking
  predictions against specific observations.
\item The gap between fiber (general model) and section (specific
  observation) is exactly where errors live, and the verification
  step systematically closes that gap.
\end{itemize}

The 33\% human-normalized score validates the pattern: even with
current-generation LLMs as the hypothesis generator, the
generate-and-verify architecture achieves qualitatively different
results (33\% vs.\ $<$1\%) compared to LLMs alone. The verification
step is what makes the difference --- it converts noisy fiber
proposals into reliable fiber representations.
\end{proposition}

%% =================================================================
\section{Cross-references}
%% =================================================================

\begin{remark}[Connections to other formalizations]
\begin{itemize}
\item \textbf{Seeding RSI} (2026-07-28): Hyperon architecture. The
  ARC-AGI-3 system is a concrete instance of Hyperon's
  neural-symbolic philosophy.
\item \textbf{Un-Jailbreakable AI} (2026-07-10): generate-and-verify.
  Same structural pattern (LLM generates, separate system verifies),
  applied to world-modeling rather than safety.
\item \textbf{Wild Times in the AGI Lab} (2026-05-26): OmegaClaw.
  Both systems use LLM + meta-control; the ARC system adds formal
  verification.
\item \textbf{The Orchard Bug} (2026-06-05): cocycle condition.
  The verifier prevents Orchard-bug-type composition failures by
  checking all transitions, not just the most recent.
\item \textbf{What Should an AGI Curriculum Look Like?} (2026-07-22):
  learning architecture. The hypothesis-revision loop is a form of
  curriculum --- the agent learns the world's rules in order of
  discoverability.
\end{itemize}
\end{remark}

\end{document}
